\begin{figure}[htbp]
  \centering
  \resizebox{0.98\textwidth}{!}{%
  \begin{tikzpicture}[
    usecase/.style={ellipse, draw=black, thick, align=center, text width=3.5cm, minimum height=1.15cm, font=\small},
    actor/.style={font=\small\bfseries, align=center},
    >=Stealth
    ]

    \coordinate (A) at (0,0);
    \draw[thick] (0, 0.8) circle (0.25);
    \draw[thick] (0, 0.55) -- (0, -0.5);
    \draw[thick] (-0.6, 0.2) -- (0.6, 0.2);
    \draw[thick] (0, -0.5) -- (-0.4, -1.2);
    \draw[thick] (0, -0.5) -- (0.4, -1.2);
    \node[actor] at (0, -1.65) {Cán bộ quản lý\\đô thị};

    \node[usecase] (uc1) at (-7, 4.4) {Thêm tài sản mới\\trên bản đồ};
    \node[usecase] (uc2) at (-7, 2.6) {Chỉnh sửa thông tin\\tài sản};
    \node[usecase] (uc3) at (-7, 0.8) {Xóa tài sản};
    \node[usecase] (uc4) at (-7, -1.0) {Tìm kiếm và lọc\\tài sản trên bản đồ};
    \node[usecase] (uc5) at (-7, -2.8) {\gls{ai} phân tích ảnh\\và phân loại tài sản};

    \node[usecase] (uc6) at (7, 4.4) {Tạo công việc\\bảo trì};
    \node[usecase] (uc7) at (7, 2.6) {Cập nhật tiến độ\\bảo trì};
    \node[usecase] (uc8) at (7, 0.8) {Theo dõi tiến độ\\bảo trì};
    \node[usecase] (uc9) at (7, -1.0) {Tra cứu lịch sử\\bảo trì tài sản};
    \node[usecase] (uc10) at (7, -2.8) {Tiếp nhận và xử lý\\phản ánh sự cố};

    \node[usecase] (uc11) at (0, 4.8) {Xem Dashboard\\tổng quan};
    \node[usecase] (uc12) at (0, 3.0) {Xem báo cáo\\thống kê tài sản};
    \node[usecase] (uc13) at (0, 1.2) {Xuất báo cáo\\và dữ liệu};

    \foreach \uc in {uc1,uc2,uc3,uc4,uc5,uc6,uc7,uc8,uc9,uc10,uc11,uc12,uc13} {
      \draw[-Stealth, thick] (A) -- (\uc);
    }

  \end{tikzpicture}%
  }
  \caption{Biểu đồ Use Case cho tác nhân Cán bộ quản lý đô thị}
  \label{fig:uc_canbo_dothi}
\end{figure}
